Prompt injection has been treated for four years as a problem of
robustifying the model against natural-language attack inputs. That
framing is structurally bounded by the open-grammar nature of the input
surface; the empirical ceiling on text-side defences against
state-of-the-art benchmarks sits at 14-31\% attack success rate.

We have argued that the consequential half of an agent's behaviour
passes through a much narrower interface, that the interface is bounded
enough to admit cryptographic and proof-based enforcement, and that the
resulting composition --- SMT-verified policy completeness over the JSON
Schema, Ed25519-signed capability tokens with mechanised attenuation
invariants, gate enforcement on the only path to tool execution ---
eliminates the tool-call-mediated prompt-injection class. The cost is
well under 100 microseconds per call on commodity hardware and the
effort of authoring a schema and a policy per tool.

The deeper claim is that the trajectory of AI security must be to move
enforcement boundaries out of the model and into structural gates
wherever the grammar permits. Tool calls are the easiest such boundary;
database queries, generated code, and structured outputs admit similar
treatment. The thesis of this paper, generalised, is that the
prompt-injection problem becomes tractable exactly when one stops trying
to solve it inside the model.

The reference implementation, Lean development, and benchmark harness
are released as open source under a permissive licence (Apache-2.0) and
available at \texttt{{[}anonymised{]}}.

\begin{center}\rule{0.5\linewidth}{0.5pt}\end{center}

\subsection{References}\label{references}

\emph{Bibliography in \texttt{paper/references.bib}. Three entries
marked \texttt{\%\%\ UNVERIFIED} need a final check against published
proceedings before camera-ready.}

\begin{itemize}
\tightlist
\item
  Anthropic. Constitutional Classifiers. 2024.
\item
  Backes, J. et al.~Semantic-based automated reasoning for AWS access
  policies using SMT (Zelkova). FMCAD 2018.
\item
  Birgisson, A. et al.~Macaroons: cookies with contextual caveats. NDSS
  2014.
\item
  Chen, S. et al.~StruQ: defending against prompt injection with
  structured queries. USENIX Sec 2024.
\item
  de Moura, L. and Bjørner, N. Z3: an efficient SMT solver. TACAS 2008.
\item
  de Moura, L. and Ullrich, S. The Lean 4 theorem prover and programming
  language. CADE 2021.
\item
  Debenedetti, E. et al.~AgentDojo: dynamic environment for attacks and
  defences for LLM agents. NeurIPS 2024.
\item
  Erbsen, A. et al.~Simple high-level code for cryptographic arithmetic,
  with proofs (Fiat-Crypto). S\&P 2019.
\item
  Fisler, K. et al.~The Margrave tool for firewall analysis. ICSE 2005.
\item
  Greshake, K. et al.~Not what you've signed up for: indirect prompt
  injection in LLM-integrated applications. AISec 2023.
\item
  Hines, K. et al.~Defending against indirect prompt injection with
  spotlighting. arXiv 2024.
\item
  Klein, G. et al.~seL4: formal verification of an OS kernel. SOSP 2009.
\item
  Levy, H. Capability-based computer systems. Digital Press 1984.
\item
  mathlib community. mathlib4: a library of mathematics for Lean 4.
  2024.
\item
  Miller, M. Robust composition: a unified approach to access control
  and concurrency control. PhD thesis, 2006.
\item
  Newman, Z. et al.~Sigstore: software signing for everybody. CCS 2022.
\item
  Perez, F. and Ribeiro, I. Ignore previous prompt: attack techniques
  for language models. arXiv 2022.
\item
  PyCA. The cryptography library. https://cryptography.io
\item
  Samuel, J. et al.~Survivable key compromise in software update
  systems. CCS 2010.
\item
  Shapiro, J. et al.~EROS: a fast capability system. SOSP 1999.
\item
  Torres-Arias, S. et al.~in-toto: farm-to-table guarantees for bits and
  bytes. USENIX Sec 2019.
\item
  Zhan, Q. et al.~InjecAgent: benchmarking indirect prompt injections in
  tool-integrated LLM agents. ACL Findings 2024.
\end{itemize}

\begin{center}\rule{0.5\linewidth}{0.5pt}\end{center}

\subsection{Author's note on remaining pre-submission
work}\label{authors-note-on-remaining-pre-submission-work}

Most \texttt{{[}TBD{]}} markers from the 2026-05-14 draft have been
resolved over the 2026-05-15 → 2026-05-17 empirical sweep (see
\texttt{paper/eval/runs/README.md} for the run history). The remaining
items before camera-ready:

\begin{enumerate}
\def\labelenumi{\arabic{enumi}.}
\tightlist
\item
  \textbf{InjecAgent LLM-driven measurement.} §6.2.1's static bound
  (0/2108) is in hand and is the structural claim. An LLM-driven run on
  InjecAgent would parallel the §6.2/§6.2.2/§6.2.3 AgentDojo evidence
  and is in flight as compute becomes available.
\item
  \textbf{A closed-frontier model row.} All measured base models are
  open-weight: DeepSeek (v3.2 / v4-flash / v4-pro, MoE), Alibaba
  (Qwen3.5:397b, MoE), Moonshot (Kimi-k2.6, MoE) --- five model
  identities from three providers, reported in §6.2. A Claude or GPT row
  would broaden the architectural diversity claim to include a
  closed-frontier provider and at least one non-MoE architecture. Gated
  on third-party credit grants (next decision cycle: 2026-06-01). We
  additionally attempted Mistral-large-3:675b (dense) and Gemma3:27b
  (Google) as substitutes; neither model is presently served on the
  Ollama Cloud tier we use for inference, so both attempts surfaced as
  \texttt{APIStatusError} rejections rather than measurements. We report
  this honestly rather than substitute a third-party-hosted alternative
  whose throughput and stochasticity we cannot equate to the rest of the
  harness.
\item
  \textbf{StruQ defence row.} StruQ requires fine-tuning the base model
  and is not a config flip in the AgentDojo harness; whether to attempt
  it before submission is a feasibility judgement.
\item
  \textbf{Anonymisation for double-blind submission --- applied.}
  Product name, module names, licence specifics, and publication dates
  softened to category-level descriptors in the body of the paper. The
  author block in \texttt{paper/main.tex} was anonymised on the day the
  LaTeX shell was created. Camera-ready de-anonymisation reverts to the
  \texttt{pre-anon} git tag's state of the affected files ---
  five-minute mechanical operation.
\item
  \textbf{Final pass for figure references and citation keys.}
  Bracket-style placeholders in the bibliography need replacing with the
  venue's chosen citation style.
\end{enumerate}

The empirical evidence base --- three generations of one open-weight
family on banking × six defences (complete), \textbf{three} generations
on three cross-suites × \{none, vcd\_full\} (19 of 21 cells; the two
omitted cells crashed in an AgentDojo slack-suite utility-scorer bug),
\textbf{three} generations on three attack families × \{none, shields,
vcd\_full\} (with full defence sweep on v4-flash), two additional
providers (Qwen3.5, Kimi-k2.6) on banking × \{none, shields,
vcd\_full\}, 8,073 logged gate decisions, 720 LLM-driven banking attack
attempts at 100\% gate-block rate outside one benchmark coincidence
cell, and full ablation across the two VCD components --- is settled.

\textbf{2026-09-01 addendum.} Since the empirical freeze, the reference
implementation gained three engineering components described in §5: the
policy-DSL gateway (network enforcement point with
per-agent/per-source/per-destination rules), remote-key signing (AWS KMS
/ Azure Key Vault / HashiCorp Vault Transit), and the conformance
re-verification layer (field-level Merkle commitments with selective
disclosure). None of these affects the measured claims: the gate, token,
and proof semantics are unchanged, and §6's numbers were produced
against the frozen pre-registration hashes. The gateway and conformance
layers are engineering contributions reported for completeness; the
paper's evaluation claims remain scoped to the gate + receipt core that
the empirical freeze covered.
